%------------------------------------------------------------------------------------
%  SECTION 5.2 : VALORISATION VISUELLE AVEC LES BOÎTES TCOLORBOX ET TIKZ
%------------------------------------------------------------------------------------
\section{Encadrés Typographiques Modernes et Schémas Vectoriels TikZ}
\label{sec:boites_tcolorbox_tikz}

L'aération d'un document passe par l'insertion rythmée d'éléments d'accroche visuelle.
En \LaTeX{} classique, la création d'encadrés personnalisés imposait des constructions complexes à base de \texttt{\textbackslash fbox} ou \texttt{mdframed}.
La classe \texttt{ClasseGuillaume} intègre le moteur de référence \textsf{tcolorbox} (avec les extensions \texttt{skins} et \texttt{breakable}), offrant des boîtes insensibles aux sauts de page et dotées d'une finition graphique irréprochable.

\subsection{L'environnement \texttt{cadreFigure} pour sublimer les schémas}
\label{subsec:cadre_figure}

L'environnement \texttt{cadreFigure} habille son contenu d'un contour raffiné aux coins délicatement arrondis (\texttt{arc = 10pt}) et d'un liseré de couleur primaire assombrie :

\begin{figure}[htbp]
  \centering
  \begin{cadreFigure}[width=0.92\textwidth]
    \begin{tikzpicture}[
      node distance = 2.1cm,
      every node/.style = {font=\small\sffamily},
      block/.style = {draw=primaryColor, thick, fill=primaryColor!8, rounded corners=3pt, minimum width=2.4cm, minimum height=1.1cm, align=center},
      sum/.style = {draw=secondaryColor, thick, circle, fill=secondaryColor!10, minimum size=0.8cm},
      arrow/.style = {->, >=stealth, thick, color=darkgray}
    ]
      % Noeuds du schéma bloc
      \node[sum] (sommateur) {$\Sigma$};
      \node[block, right of=sommateur, xshift=1.2cm] (correcteur) {Correcteur\\Numérique $C(z)$};
      \node[block, right of=correcteur, xshift=1.6cm] (processus) {Procédé Physique\\Modélisé $G(s)$};
      \node[coordinate, right of=processus, xshift=1.2cm] (sortie) {};
      \node[coordinate, below of=correcteur, yshift=0.8cm] (retour) {};

      % Flèches de connexion
      \draw[arrow] ([xshift=-1.6cm]sommateur.west) -- node[above] {$r(t)$} node[below, pos=0.2] {\footnotesize Consigne} (sommateur.west);
      \draw[arrow] (sommateur) -- node[above] {$e(t)$} (correcteur);
      \draw[arrow] (correcteur) -- node[above] {$u(t)$} (processus);
      \draw[arrow] (processus) -- node[above] {$y(t)$} (sortie);
      \draw[arrow] ([xshift=0.8cm]processus.east) |- (retour) -| node[pos=0.9, left] {$-$} (sommateur.south);
    \end{tikzpicture}
  \end{cadreFigure}
  \caption{Architecture fonctionnelle d'une boucle d'asservissement en temps réel composée sous TikZ}
  \label{fig:schema_bloc_tikz}
\end{figure}

Le schéma vectoriel représenté à la~\cref{fig:schema_bloc_tikz} illustre la symbiose parfaite entre le graphisme vectoriel natif \textsf{TikZ} et le conteneur \texttt{cadreFigure}.

\subsection{L'environnement \texttt{boiteSynthese} pour les messages clés}
\label{subsec:boite_synthese}

L'environnement \texttt{boiteSynthese\{Titre\}} structure les conclusions d'étapes, définitions ou théorèmes au moyen d'un fond teinté doux (\texttt{primaryColor!6}) et d'un titre en lettres sans-serif grasses :

\begin{boiteSynthese}{Principe Fondamental de l'Équilibre Graphique}
  Les boîtes décoratives doivent soutenir la transmission intellectuelle et non la distraire :
  \begin{itemize}
    \item Une boîte ne doit contenir qu'une seule idée forte ou synthèse transversale.
    \item L'intensité de la couleur d'arrière-plan ne doit jamais dépasser \qty{8}{\percent} pour garantir un ratio de contraste d'impression optimal et éviter la fatigue oculaire.
    \item Dans un mémoire volumineux, limitez l'usage à une boîte de synthèse par section majeure.
  \end{itemize}
\end{boiteSynthese}

\subsection{Boîtes d'alerte avec bandeau vertical (\emph{Left-bar alert})}
\label{subsec:boite_alerte}

Pour les mises en garde méthodologiques ou les avertissements de sécurité expérimentale, un style minimaliste avec filet latéral (\emph{leftbar}) s'avère particulièrement sobre et percutant :

\begin{tcolorbox}[
  enhanced,
  blanker,
  borderline west={3pt}{0pt}{Bittersweet},
  left=12pt,
  top=4pt,
  bottom=4pt
]
  \textcolor{Bittersweet}{\bfseries\sffamily Règle de Sécurité et Hygiène des Données :}\\
  \small\itshape Lors de l'acquisition de signaux physiologiques sur sujets humains, assurez-vous impérativement de la conformité avec le Règlement Général sur la Protection des Données (RGPD) et de l'anonymisation irréversible des identifiants patients avant tout traitement statistique.
\end{tcolorbox}

